\documentclass[a11paper, 11pt]{article}

\usepackage{libs/document}
\usepackage{libs/titlepage}
\usepackage{libs/codeblocks}
\usepackage[T1]{fontenc}
\usepackage[french]{babel}

% \addbibresource{bibliography.bib}
% \nofiles

\newcommand{\todo}[1]{\textcolor{orange}{\textbf{TODO}: #1}}
\newcommand{\note}[1]{\textcolor{purple}{\textbf{NOTE}: #1}}
\newcommand{\xxx}[1]{\textcolor{red}{\textbf{XXX}: #1}}

% \institution{Université de Sherbrooke}
% \faculty{Faculté de génie}
% \department{Département de génie électrique et de génie informatique}
\title{Rapport d'APP}
\classnb{GIF391}
\class{Conception d'un système distribué}
\author{
  \addtolength{\tabcolsep}{-0.4em}
  \begin{tabular}{rcl} % Ajouter des auteurs au besoin
  Benjamin Chausse & -- & CHAB1704 \\
  Samuel Bilodeau  & -- & BILS2704 \\
  \end{tabular}
}
\teacher{Frédéric Mailhot}
% \location{Sherbrooke}
% \date{\today}

\begin{document}
\maketitle
\newpage
\tableofcontents
\newpage

\section{Description des solutions utilisées}

\subsection{Arrêt 1: Exemple d'utilisation des commandes}

À ce stade, l'infrastructure emploie un seul conteneur pour son fonctionnement.
Afin de lancer le système, il faut d'abord construire l'image du conteneur à
partir de son \textit{Dockerfile}. Par la suite, il faut lancer le conteneur
afin qu'il puisse s'exécuter en arrière-plan. Enfin, il est possible de s'y
connecter pour utiliser le système. Voici un exemple d'utilisation de
l'infrastrucure (le résultat de commandes de mise en place est remplacé par
"\texttt{...}"):

\begin{code}[title={zsh (dans un terminal)},listing options={language=bash, style=bashstyle}]
docker build -t arret01 . # Construction de l'image
docker run -it arret01 # Lancement interactif du conteneur
\end{code}

\begin{code}[title={sh (dans le conteneur)},listing options={language=bash, style=bashstyle}]
liste   # Liste des taches dans la liste (devrait etre vide)
nouveau # Ajout d'une taches utilisant fortune pour la description
liste   # La liste devrait maintenant contenir un taches
retrait # Retrait de la tache ajoute precedemment
liste   # La liste devrait etre vide a nouveau
\end{code}

À l'intérieur de l'environnement de l'application (dans le conteneur), les
commandes \texttt{liste}, \texttt{nouveau} et \texttt{retrait} permettent la
gestion des tâches pour un seul utilisateur. Puisque ce système sert de
prototype, la commande \texttt{nouveau} génère une description aléatoire à
l'aide de l'utilitarie \texttt{fortune}. Dans un système réel, cette commande
devrait être remplacée par une commande permettant d'ajouter une tâche avec une
description spécifiée par l'utilisateur et l'auteur de la tâche devrait être
déterminé au préalable (possiblement dans un fichier de configuration).

\subsection{Arrêt 2: Fichier de configuration fonctionnel}
\input{arret02.tex}

\newpage

Lors de cet arrêt, l'application est segmentée en deux section. Le client ne
fait qu'envoyer des requêtes mais ne contient aucune tâche. Le serveur lui,
gère les tâches et répond aux requêtes du client. Ce type de segmentation peut
permettre en pratique d'avoir le serveur sur une machine et le client sur une
autre. Dans le cas de la problématique toutefois, le client et le serveur sont
sur la même machine afin de simplifier les tests.

% \newpage
\subsection{Arrêt 3: Fichier \textit{yml} utilisé pour \textit{docker-compose}}
\input{arret03.tex}

Dans cet arrêt, l'idée est la même qu'à l'arrêt 2, mais cette fois-ci, trois
clients utilisent le service (chacun ayant un serveur unique leur étant
associé). L'avantage de cette approche est (en plus des mêmes présentes à
l'arrêt 2) qu'aucun de conflit de données ne peut survenir entre les clients
puisque leurs données sont stockées sur des serveurs différents.

\subsection{Arrêt 4: Explication des problèmes inhérents}

Cet arrêt a pour but de simplifier l'infrastructure de l'arrêt précédent en
implémentant un seul serveur pour tous les clients. Cette approche est plus
simple si elle est bien mise en place. Toutefois, l'implémentation actuelle est
vulnérable à des confilts de données. Des problèmes peuvent survenir si deux
clients (ou plus) tentent d'accéder aux tâches en même temps. Ceci est dû au
fait que toutes les tâches sont stockées dans un seul fichier.

\subsection{Arrêt 5: Fichier \textit{yml} utilisé pour \textit{docker-compose}}
\input{arret05.tex}

Compte tenu des problèmes de l'arrêt précédent, l'objectif de cet arrêt est de
rendre l'infrastructure plus robuste. Pour ce faire, les données sont stockées
dans une base de donnée \textit{PostgreSQL}. Cette base de donnée s'occupe de
la gestion des conflits de données.


\subsection{Arrêt 6: Fichier \textit{yml} utilisé pour \textit{kubernetes}}

Lors de ce dernier arrêt, le but but est de simuler un déploiement de
l'infrastructure sur plusieurs machines. Pour simuler ce type de déploiement,
\textit{minikube} est utilisé. \textit{Minikube} est un outil qui permet à
\textit{kubernetes} (un outil permettant de gérer des conteneurs distribués).

\input{arret06.tex}


\section{Discussion de la structure}

La version finale du système offre une plus grande flexibilité par rapport aux
premières itérations. Son principal inconvénient réside dans sa complexité
croissante, ce qui peut entraîner des problèmes à différents niveaux, parfois
difficiles à détecter. Cependant, les avantages sont nombreux. La nature
distribuée du système facilite son expansion horizontale, le rendant ainsi
adapté à un grand nombre d'utilisateurs. L'utilisation d'une base de données
garantit l'atomicité des transactions pour éviter les conflits et les
enregistrements en double. La réplication des serveurs et des bases de données
augmente la redondance et, par conséquent, la résilience du système. La
modularité des composants permet des mises à niveau fluides et transparentes
pour l'utilisateur.

\subsection{Technologies sous-jacentes}

\subsubsection{Identification des ressources}

La technologie utilisée pour identifier les ressources s'appelle "namespace".
Elle nous a permis d'attribuer différentes ressources. Au début, nous l'avons
principalement utilisée pour assigner des emplacements de stockage afin de
partager des fichiers entre les différents conteneurs. Par la suite, lorsque
des requêtes réseau via SSH ont été introduites, nous avons utilisé les
namespaces pour créer un réseau virtuel sur lequel les machines pouvaient
communiquer facilement. Les noms des conteneurs étaient utilisés pour remplacer
les adresses IP, ce qui simplifiait grandement nos interactions. En résumé, les
namespaces sont utiles pour assigner diverses ressources telles que le réseau,
les systèmes de fichiers, les processus, les groupes ou les utilisateurs, et
bien d'autres encore.

\subsubsection{Contrôle d'accès aux ressources}

Pour contrôler l'accès aux ressources, nous avons utilisé les cgroups. Ces
derniers permettent d'attribuer l'allocation de différentes ressources
(processeur, mémoire, disque) en fonction des membres du cgroup. Ces membres
peuvent être un processus spécifique, un utilisateur ou un conteneur dans son
intégralité. Nous pouvons gérer la quantité minimale de ressources et donner la
priorité à certaines entités par rapport à d'autres. Les cgroups sont également
un bon moyen de gérer les accès à des répertoires et à des réseaux spécifiques.

\subsubsection{Gestion des accès aux fichiers utilisés}


Docker utilise une technologie Linux très pratique appelée "système de fichiers
en couches" (Union File System). Le concept est simple : certains répertoires
sont en lecture seule, et lorsqu'une modification est nécessaire, elle est
enregistrée dans un répertoire qui stocke uniquement les différences (ce qui
est conceptuellement similaire la technologie Git). Cette approche permet à
plusieurs conteneurs d'utiliser la même base commune, mais chacun d'entre eux
de manière distincte avec son propre répertoire de modifications. Dans la
dernière itération, plusieurs clients utilisaient l'image client, mais chacun
était unique. Cela nous permet d'utiliser efficacement l'espace de stockage et
d'éviter de corrompre l'image initiale.


% \todo{Quels sont les pilotes utilisables pour la persistance du système de fichiers
% à union ?} \\
% \todo{Quel pilote de persistence a été utilisé et pourquoi ?}

Plusieurs pilotes sont disponibles pour gérer les fichiers à union, tels que
UFS (Advanced Multi-Layered Unification File System), DeviceMapper, Btrfs
(B-Tree File System), ZFZ (Zettabyte file system), mais celui qui nous
intéresse est OverlayFS, plus spécifiquement sa deuxième version, Overlay2.
Depuis la version 1.13 de Docker, Overlay2 est le pilote par défaut utilisé.

Overlay2 est une version améliorée d'OverlayFS qui nous permet de gérer
efficacement les modifications tout en maintenant la transparence des fichiers
uniques pour l'utilisateur. Il garantit également l'intégrité des fichiers
sous-jacents et réduit l'utilisation de l'espace mémoire nécessaire. Avec
Overlay2, les modifications apportées aux fichiers dans les conteneurs sont
enregistrées de manière efficace en créant des couches supplémentaires, tout en
préservant l'état d'origine de l'image de base. Cela permet une utilisation
efficace des ressources et facilite la création de conteneurs basés sur des
images existantes.

\subsubsection{Configuration réseau pour la communication des conteneurs}

Pour simplifier l'environnement dans lequel nous travaillions, nous avons
utilisé un namespace de réseau pour regrouper toutes les entités sur la même
couche. Bien que cela fonctionnait, cette architecture n'était pas la plus
idéale dans un environnement de production. Certaines entités, comme les
clients, ne devraient pas nécessairement avoir accès aux autres, telles que la
base de données. En général, ceux-ci sont disposées sur des réseaux distants et
les interactions requises sont déterminées à l'avance. Afin d'assurer
l'encapsulation des composants et la sécurité du système, l'authentification et
le chiffrement sont des étapes obligatoires avant toute action. Étant donné que
l'objectif de l'activité ne portait pas spécifiquement sur la mise en réseau,
tous les services étaient visibles et accessibles à tous.

\subsubsection{Duplication mise en place pour les ressources}

Il y a quatre types principaux de ressources dans notre système : matériel,
logiciel, données et réseau. La duplication de ces ressources s'applique
différemment en fonction du contexte. En ce qui concerne le matériel, nous
avons développé nos solutions dans un environnement de test principalement sur
nos ordinateurs portables et systèmes personnels. Toutes les composantes
utilisaient des ressources virtuellement dupliquées, mais fondamentalement,
nous utilisions le même processeur et la même mémoire vive pour faire
fonctionner l'ensemble du système. Pour le logiciel, toutes les images
reposaient sur les mêmes fondations Docker et contenaient une copie du script
"todo.py". Cela signifiait que toute modification nécessitait la mise à jour de
tous les fichiers et la relance complète du déploiement. En ce qui concerne les
données, nous sommes passés d'un simple fichier texte à une base de données
entièrement autonome. Dans le cas où plusieurs bases de données seraient
utilisées, la réplication des données entre les bases de données deviendrait un
aspect important à prendre en compte. En ce qui concerne le réseau, étant donné
que nous étions tous sur la même machine, aucune duplication n'était
nécessaire. Il est important de noter que la duplication des ressources peut
varier en fonction des exigences spécifiques du système pour lequel il est
conçu, de la redondance souhaitée et des contraintes de performance. Les
stratégies de duplication sont mises en place pour assurer la disponibilité et
la résilience des opérations du système.


% \todo{Quelles ressources doivent être dupliquées, et pourquoi?} \\

Dans le cas d'un système ou tout les clients/serveurs sont sur la même machine
(ex: arrêt 3), la duplication d'images est inutile puisque chaque
serveur/client peuvent utiliser les mêmes images. Toutefois, dans le cas d'un
système distribué (ex: si l'arret 6 était réparti sur plusieurs machines), il
serait nécéssaire que chaque image soit dupliquée sur chaque machine. En ce qui
concerne les données, la duplication dépends de la redondance souhaitée.


Dans un système de production, la duplication des données est essentielle pour
garantir leur intégrité. On met généralement en place des répliques de bases de
données et des copies de sauvegarde qui sont envoyées en dehors de leur site
physique respectif. Cette duplication des données améliore également la
résilience et les performances du système. En répartissant les sites
géographiquement, on réduit la distance entre l'utilisateur et le système, et
en cas de panne, le trafic peut être redirigé vers d'autres sites fonctionnels.

De manière générale, les systèmes distribués sont situés dans des clusters où
la duplication est déjà prise en charge. Les machines physiques ne sont que des
unités de puissance de calcul, et la défaillance de l'une d'entre elles
n'entraîne pas de perte de performance notable. La redondance du réseau est
également assurée dans ces centres, car ils disposent d'une infrastructure
robuste face aux pannes. De nos jours, les systèmes de cloud se chargent en
grande partie de la réplication physique pour la performance et de la
redondance réseau pour assurer l'accessibilité.

Dans un environnement de développement ou de test, il est courant de dupliquer
les données uniquement sur notre machine pour faciliter les opérations. En
revanche, dans un système de production, la duplication des données devrait se
faire sur plusieurs sites pour plusieurs raisons. Tout d'abord, la
disponibilité : en cas de panne d'une machine, une autre machine dans un autre
site peut prendre en charge rapidement les demandes. Ensuite, la performance :
avec des sites géographiquement dispersés, les données sont réparties et les
requêtes sont acheminées vers le serveur le plus proche, réduisant ainsi le
temps de réponse.

De façon connexe, il est possible de dupliquer les données d'un service dans
une même machine. Prenons l'exemple suivant: un service n'est pas suffisamment
grand pour justifier l'utilisation de plusieurs machines qui hébergent des
données. Ces données sont stockés dans des disques durs ayant une période de
vie finie. Dans le cas où un disque dur tombe en panne, les données sont
perdues. Si cette information est dupliquée sur un autre disque, la perte de
données est évitée sans avoir à utiliser plusieurs machines.

\section{Conclusion}

Pour conclure, les systèmes distribués font partie intégrante de notre
quotidien, car ils sont présents dans toutes les applications et services que
nous utilisons. Leur objectif principal est d'assurer l'intégrité, la
confidentialité et la disponibilité des données. Ils parviennent à atteindre
cet objectif grâce à la duplication des ressources aux quatre niveaux
mentionnés : matériel, logiciel, données et réseau.

La redondance au niveau matériel et réseau garantit l'accessibilité du système,
en permettant la reprise rapide en cas de défaillance d'une machine ou d'un
réseau. La duplication des composants logiciels et des données contribue à
améliorer les performances et la disponibilité en répartissant les charges de
travail et en permettant une récupération rapide en cas de panne.

Les systèmes distribués et leur duplication des ressources jouent un rôle
crucial dans notre monde connecté. Ils nous permettent de bénéficier
d'applications et de services fiables et disponibles en permanence, améliorant
ainsi notre expérience numérique du quotidien.

% \newpage
% \printbibliography[heading=bibintoc]
\end{document}
